\documentclass[10pt,a4paper]{article}
\usepackage[T1]{fontenc}
\usepackage[utf8]{inputenc}
\usepackage[english]{babel}
\usepackage{amsmath,booktabs,float,geometry,graphicx,hyperref}
\geometry{margin=1.6cm}
\hypersetup{colorlinks=true,urlcolor=blue,linkcolor=blue}

\title{Exact Product Navigation for Hard-to-Ground PDDL VisitAll}
\author{Technical Report on \texttt{pddl\_fol}}
\date{September 2026}

\begin{document}
\maketitle

\begin{abstract}
The lifted \texttt{folplan} planner can avoid Cartesian action grounding, but
its generic goal-cover estimate provides little direction on VisitAll tasks
with few target tuples. We identify a structural fragment in which movement
changes one coordinate of a location tuple along a static neighbor graph and
marks each reached tuple as visited. For up to three remaining targets, we
derive exact directed distances and construct a shortest plan from the best
target order. The optimizer recognizes action structure rather than domain
names and replays every constructed plan against the original PDDL semantics.
On 54 preregistered hard-to-ground VisitAll instances, it solves all 54 within
15 seconds and 1 GiB; Powerlifted and Tyr each solve 35. On their 35 shared
solved instances, the median paired full-process speedups are 87.5 and 16.7,
respectively. These results establish competitiveness on this narrow
structural fragment, not on general PDDL planning.
\end{abstract}

\section{Motivation}
The hard-to-ground (HTG) VisitAll corpus has 180 instances in 18 families:
three, four, or five location coordinates; close or far target placement;
and one to three goal tuples~\cite{htg,wichlacz2022}. The action schemas have
one parameter for every old coordinate and one for the new coordinate, so
complete grounding grows rapidly with object count. Lifted joins already
avoid this grounding cost in \texttt{folplan}~\cite{previous}.

Our earlier eight-instance study used only the five-dimensional close,
single-target family~\cite{previous}. Every unsolved goal in that selection
has exactly one missing \texttt{visited} atom. Consequently, the goal-cover
heuristic is one in every non-goal state; the full Cartesian action universe
also exceeds the threshold for delete-relaxed $h_{\max}$. On \texttt{p9},
the target differs from the initial position by one, seven, and five edges
in three coordinates, respectively. Its exact travel distance is 13, while
the generic estimate remains one. The previous A* run reached the 1 GiB
memory limit on this task.

\section{Recognized fragment and exact reduction}
Recognition examines the validated task's predicate and action structure.
Exactly one dynamic location tuple is true initially. Each action requires
that tuple and one static directed neighbor fact, changes exactly one tuple
coordinate, deletes the old location, adds the new location, and adds a
\texttt{visited} fact for the new location. No action deletes \texttt{visited},
and the goal is a positive conjunction of at most three ground
\texttt{visited} tuples not already true. Task, predicate, parameter, and
action names play no role. Tasks outside this fragment use generic A*.

Let $q=(q_1,\ldots,q_m)$ be a location and let $d_N(u,v)$ be shortest-path
distance in the directed static neighbor graph. Independent coordinate moves
give an exact location distance
\begin{equation}
  d(q,r)=\sum_{j=1}^{m}d_N(q_j,r_j).
  \label{eq:product-distance}
\end{equation}
An unreachable coordinate pair makes the corresponding distance infinite.
Let $R$ contain remaining target tuples and let $q_0$ be the initial
location. The optimal route cost is
\begin{equation}
  C^*=\min_{\pi\in\mathrm{Perm}(R)}
       \sum_{i=0}^{|R|-1} d(q_i,q_{i+1}),
  \quad q_{i+1}=\pi_{i+1}.
  \label{eq:route}
\end{equation}
For at most three targets, this considers at most six orders. Any valid plan
has an order in which it first reaches the targets; each segment costs at
least its distance in Equation~\ref{eq:product-distance}. Conversely,
concatenating shortest coordinate paths for a minimizing order visits every
target at that cost, possibly visiting some targets earlier along a segment.
Thus Equation~\ref{eq:route} is exact. The constructed ground actions are
replayed with the ordinary PDDL evaluator before the plan is returned.
Explicit BFS retains its generic graph-search implementation.

\section{Evaluation protocol}
We preregister \texttt{p0}, \texttt{p5}, and \texttt{p9} from each of the 18
official VisitAll families at corpus commit \texttt{c5670ed2}: 54 tasks,
including easy, middle, and deep instances in every dimension, distance,
and goal-count group. The independent metric oracle computes expected
optimal lengths for all 180 official tasks; the selected tasks are checked
against that oracle and the returned plans against \texttt{folplan}'s exact
replay.

We compare release builds of \texttt{folplan}, Powerlifted's
\texttt{alt-bfws1} with FF and Yannakakis joins~\cite{powerlifted}, and
Tyr's single-worker lifted greedy best-first search with RPG-FF~\cite{tyr}.
Each process has a 1 GiB address-space cap and a 15-second wall timeout.
Solved cases receive three interleaved runs; failures are attempted once.
Reported time includes process startup, parsing or translation, and search.
Plan lengths and outcome categories accompany times. Compiler flags, exact
revisions, executable hashes, raw rows, and commands are archived beside
the benchmark harness.
Primary endpoints are task coverage and paired median full-process latency
against Tyr; Powerlifted gives a second lifted reference. We report the
number of tasks won by each solver as well as the distribution of paired
runtime ratios, retaining failed tasks in the coverage denominator.

\section{Results}
Table~\ref{tab:coverage} gives complete coverage and failure categories. All
successful solver--task cells completed three runs; no cell had a later-run
failure. Powerlifted's 16 memory-limit results are exit-code 22 with explicit
memory-limit messages in the logs; its other three failures and Tyr's 19
failures are 15-second process timeouts. The specialized planner's measured
time ranges from 0.0027 to 0.0089 seconds across tasks, including parsing and
plan replay.

\begin{table}[H]
\centering
\caption{Coverage under 15 seconds and 1 GiB per process. Median time uses
only tasks solved by that planner; paired speedups below use identical tasks.}
\label{tab:coverage}
\begin{tabular}{lrrrrrr}
\toprule
Planner & Solved & Close & Far & Timeout & Memory & Median time (s) \\
\midrule
\texttt{folplan} A* & 54/54 & 27/27 & 27/27 & 0 & 0 & 0.0037 \\
Powerlifted & 35/54 & 26/27 & 9/27 & 3 & 16 & 0.2746 \\
Tyr & 35/54 & 26/27 & 9/27 & 19 & 0 & 0.0588 \\
\bottomrule
\end{tabular}
\end{table}

On all 35 tasks solved by both \texttt{folplan} and each comparator,
\texttt{folplan} has lower median full-process wall time. The median paired
ratio $t_{\mathrm{Powerlifted}}/t_{\mathrm{folplan}}$ is 87.5
(interquartile range 41.4--530.8); the corresponding Tyr ratio is 16.7
(8.7--117.5). These ratios include startup and translation costs, so they
measure end-to-end use rather than search alone. Figure~\ref{fig:competition}
shows the complete coverage/time curve and stratified coverage. A larger
per-instance status chart accompanies the archived data.

\begin{figure}[H]
\centering
\includegraphics[height=0.78\textheight,width=\linewidth,keepaspectratio]{docs/figures/product-visitall-compact.png}
\caption{Performance on preregistered VisitAll instances. Cactus curves sort
each planner's solved tasks by its median full-process time; shorter curves
mean lower coverage. Each dimension group contains 18 tasks.}
\label{fig:competition}
\end{figure}

The specialized planner's returned length matches the archived independent
route-length oracle on all 54 measured tasks; every returned plan also passes
the planner's exact PDDL replay. Powerlifted and Tyr each return lengths equal
to that oracle on 31 of their 35 solved tasks. Their remaining four lengths
are longer; their plans were not independently replayed. The previously
memory-limited five-dimensional close \texttt{g1/p9} case now returns the
13-step optimal plan in 0.0036 seconds median full-process time.

The harness, preregistered input hashes, all 410 attempt rows, 162 task-level
summaries, failure classifications, executable hashes, and the 180-task
oracle are archived in \path{benchmarks/results/product-visitall/}. Its
commands and pinned solver revisions are recorded in \path{metadata.json}
and \path{raw.csv}; the plotting script regenerates both figures from the
archived CSVs.

\section{Scope}
This experiment tests a deliberately narrow PDDL fragment with at most
three remaining visit goals. The structural proof and replay check are
essential: applying the route formula to actions with extra preconditions,
diagonal or teleport moves, or non-monotone visits would be unsound. Full
process latency includes different startup and translation costs for the
three implementations. The 54 selected tasks are a fixed subset of the
180-instance VisitAll corpus, not a measurement of general planner quality.

\small
\begin{thebibliography}{9}
\raggedright
\bibitem{htg} A.~B. Corr\^ea, \emph{htg-domains}, corpus commit
\texttt{c5670ed2}.
\url{https://github.com/abcorrea/htg-domains/tree/c5670ed289f58d3d6e104d7b06f311cf2f87b354}.
\bibitem{wichlacz2022} J. Wichlacz, D. H\"oller, and J. Hoffmann,
``Landmark Heuristics for Lifted Classical Planning,'' IJCAI, 2022.
\url{https://www.ijcai.org/proceedings/2022/0647.pdf}.
\bibitem{powerlifted} A.~B. Corr\^ea and J. Seipp,
``Best-First Width Search for Lifted Classical Planning,'' ICAPS, 2022.
\url{https://doi.org/10.1609/icaps.v32i1.19780}.
\bibitem{tyr} D. Drexler, O. Joergensen, and J. Seipp,
``Parallel Lifted Planning via Semi-Naive Datalog Evaluation,'' arXiv:2605.07584,
2026. \url{https://arxiv.org/abs/2605.07584}.
\bibitem{previous} \texttt{pddl\_fol},
``Lifted Successor Joins and Admissible A* in a First-Order PDDL Planner,''
technical report, 2026. \path{docs/lifted-visitall-study.tex}.
\end{thebibliography}
\end{document}
